% =====================================================================
%   附录 D：知识图谱 (Knowledge Graph)
% =====================================================================

\section{知识图谱}\label{app:kg}

\subsection{数据集细节}

\subsubsection*{构造细节}

为评估 LLM 的决策能力，尤其是其长期规划水平，本文精心整理了源自
F REEBASE 上既有知识库问答（KBQA）数据集的数据，包括 GrailQA
\citep{gu2021beyond}、ComplexWebQuestions \citep{talmor2018complexwebquestions}
与 GraphQuestions \citep{su2016}。

我们将 KBQA 视为一种工具学习（tool learning）场景，为 LLM 提供一系列
KG 查询工具。借助 \citet{gu2022arcaneqa} 所标注的 S-expressions，我们能够
准确建立与每个问题相对应的最优工具调用序列。为保证任务具有足够的难度，
我们仅保留那些至少需要 5 次工具调用的问题。通过这一严格筛选流程，我们
最终累积了 1\,663 个问题的数据集。数据集中每条样本具有以下字段：
\begin{itemize}
    \item \textbf{输入问题（Input Question）}：涉及复杂 KG 信息检索的自然
          语言陈述。
    \item \textbf{主题实体（Topic Entities）}：输入问题中提到的主题实体集合。
          我们省略实体链接过程，使 LLM 能够专注于长期规划。
    \item \textbf{动作序列（Action Sequence）}：能引导得到目标答案的"金标准"
          动作序列（即工具调用序列）。
    \item \textbf{金标准答案（Gold Answer）}：问题的金标准答案，通常以一组
          KG 实体的形式呈现。
\end{itemize}

需要注意的是，相较于 AgentBench 中与数据库的交互——数据库的具体信息与内容
会被整合到输入中——向 LLM 详尽描述一个庞大的 KG 是不现实的。该任务的本质
特征即是部分可观察的环境，这也是其不可或缺的核心要素。

\subsubsection*{评测设置}

为支持评测，我们首先使用 Virtuoso\footnote{\url{https://github.com/dki-lab/Freebase-Setup}}
托管最新版本的 F REEBASE。鉴于 SPARQL 查询的复杂性，我们决定不让 LLM
自行撰写 SPARQL 查询；相反，我们实现了一系列与 Virtuoso 后端对接的 API，
从而使 LLM 能更便捷地对 KG 进行查询。

我们使用数据集中最新的 500 个任务进行评估。当任务成功执行时，理想情况下应
经历以下阶段：
\begin{enumerate}
    \item \textbf{初始化}：向 LLM 给出具体任务描述，包括对每个所提供
          KG 查询工具的具体说明。
    \item \textbf{交互}：在此阶段，LLM 被期望调用不同工具以访问 KG，并
          累积回答问题所需的信息。重要的是，整个过程是完全自主的，
          LLM 自行决定完整的工作流。
    \item \textbf{最终答案预测}：在与 KG 交互过程中，LLM 可能生成一个
          变量列表，其中每个变量代表一组独立的实体。若 LLM 确定某个
          特定的变量应当作为最终答案，则其将输出该变量并终止任务。
\end{enumerate}

\subsubsection*{评测指标}

本文以 F1 分数作为研究的主要评测指标，通过将模型的预测答案与金标准答案
进行比对来计算。除 F1 分数外，我们还采用了精确匹配（Exact Match）指标。
与先前研究基于逻辑形式衡量精确匹配不同，本文基于预测答案集合与金标准
答案集合的精确匹配来评估。

最后，我们还评估了模型所生成动作序列的可执行性（Executability）。若模型
生成的某个动作序列被执行后能产生任一答案集合，则其可执行性得分为 1.0；
若未能产生任何答案，则为 0。

\subsection{提示词示例}

\textbf{任务描述}如下：
\begin{quote}\small\ttfamily
User: You are an agent that answers questions based on the knowledge
stored in a knowledge base. To achieve this, you can use the following
tools to query the KB.\\
\\
1. \code{get_relations(variable: var) -> list of relations}\\
A variable can be either an entity or a set of entities (i.e., the result
of a previous query). This function helps to navigate all relations in
the KB connected to the variable, so you can decide which relation is
the most useful to find the answer to the question. A simple use case can
be \code{get\_relations(Barack Obama)}, which finds all relations/edges
starting from the entity Barack Obama. The argument of \code{get\_relations}
should always be an entity or a variable (e.g., \code{\#0}) and not
anything else.\\
\\
2. \code{get_neighbors(variable: var, relation: str) -> variable}\\
Given a variable, this function returns all entities connected to the
variable via the given relation. Note that, \code{get_neighbors()} can
only be used after \code{get_relations()} is used to find a set of viable
relations. A simple use case can be \code{get_neighbors(Barack Obama,
people.person.profession)}, which returns the profession of Obama in
Freebase.\\
\\
3. \code{intersection(variable1: var, variable2: var) -> variable}\\
Given two variables, this function returns the intersection of the two
variables. The two variables MUST be of the same type!\\
\\
4. \code{get_attributes(variable: var) -> list of attributes}\\
This function helps to find all numerical attributes of the variable.
Please only use it if the question seeks for a superlative accumulation
(i.e., argmax or argmin).\\
\\
5. \code{argmax(variable: var, attribute: str) -> variable}\\
Given a variable, this function returns the entity with the maximum value
of the given attribute. It can only be used after \code{get_attributes()}
is used to find a set of viable attributes. A simple use case can be
\code{argmax(variable, age)}, which returns the oldest entity belonging
to the variable.\\
\\
6. \code{argmin(variable: var, attribute: str) -> variable}\\
Given a variable, this function returns the entity with the minimum value
of the given attribute. It can only be used after \code{get_attributes()}
is used to find a set of viable attributes. A simple use case can be
\code{argmin(variable, age)}, which returns the youngest entity belonging
to the variable.\\
\\
7. \code{count(variable: var) -> int}\\
Given a variable, this function returns the number of entities belonging
to the variable.\\
\\
After a variable is produced along the process, you need to judge whether
a variable is the final answer to the question. Each variable is
represented as an id starting from 0. For example, \code{\#0} is the first
variable, \code{\#1} is the second variable, and so on. Once you find the
answer, respond with \code{`Final Answer: \#id'}, where \code{id} is the
id of the variable that you think is the final answer. For example, if
you think \code{\#3} is the final answer, you MUST respond with
\code{`Final Answer: \#3'}. You can only take ONE action at a time!! After
you get the observation from its execution, you can take another action.
You can take at most 15 actions to find the answer to the question.
\end{quote}

鉴于让 LLM 直接查询知识库本身存在固有难度，我们观察到在零样本（zero-shot）
设置下，LLM 难以生成任何具有实质意义的输出。因此，我们在提示词中额外
提供一个教学示例，示例如下：

\begin{quote}\small\ttfamily
Question: the bipropellant rocket engine with kerosene and gas-generator
cycle is designed by who?\\
Entities: [Gas-generator cycle, Kerosene]\\
\\
Thought: I need to first find engines with gas-generator cycle. To achieve
this, I will query the KB to find relations connected to the entity
`Gas-generator cycle' and see if any of them can help me find the answer.\\
Action: \code{get_relations(Gas-generator cycle)}\\
Observation: [\dots]\\
Thought: From the above relations, I might use
\code{spaceflight.rocket_engine_cycle.rocket_engines} to find the engines
of Gas-generator cycle.\\
Action: \code{get_neighbors(Gas-generator cycle, spaceflight.rocket_engine_cycle.rocket_engines)}\\
Observation: variable \#0, which are instances of
\code{spaceflight.bipropellant_rocket_engine}.\\
\\
（后续示例省略，详见原论文。）
\end{quote}
